\scanpage{87}
- \mspage{IV-1008} -
\[
\psi_M:(\operatorname{gr}_\bullet(M)\otimes_A(A/\underline J))[T_1,\ldots,T_n]
 \longrightarrow\operatorname{gr}'_\bullet(M)
\]
\[
\psi_R:(\operatorname{gr}_\bullet(R)\otimes_A(A/\underline J))[T_1,\ldots,T_n]
 \longrightarrow\operatorname{gr}''_\bullet(R)
\]
\margin{\unsure{identifier aux précédents} sans
\unsure{oublier que comme d'ident des deux dernières flèches}
verticales de 19.6.1.1.}
Comme la \deleted{\unsure{gr}} filtration $(R''_k)$ est plus fine que $(R'_k)$ sur $R$, on a un
homomorphisme canonique \inserted{$\gamma:$}
$\operatorname{gr}''_\bullet(R)\longrightarrow\operatorname{gr}'_\bullet(R)$;
nous désignerons par $\psi'_R$
l'homomorphisme composé $\gamma\circ\psi_R$. Il résulte alors aussitôt des définitions
que l'on a un diagramme commutatif

\margin{orientation du diagramme \unsure{vicieuse}; à harmoniser avec 19.6.1.1.}
(19.6.2.\unsure{1})
\diagram{\begin{tikzcd}[ampersand replacement=\&,column sep=large,row sep=large,
execute at end picture={\draw (callout) to[bend left=15] (jl);
\draw (callout) to[bend left=5] (pl);}]
\&0\arrow[d]\\
(\operatorname{gr}_\bullet(R)\otimes_A(A/\underline J))[T_1,\ldots,T_n]
 \arrow[r,"\psi'_R"]
 \arrow[d,"{\operatorname{gr}(j)\otimes1}"'{name=jl,draw,ellipse}]
\&\operatorname{gr}'_\bullet(R)\arrow[d,"\operatorname{gr}'(j)"]\\
(\operatorname{gr}_\bullet(\unsure{N|M})\otimes_A(A/\underline J))[T_1,\ldots,T_n]
 \arrow[r,"\psi_N"]
 \arrow[d,"{\operatorname{gr}(p)\otimes1}"'{name=pl,draw,ellipse}]
\&\operatorname{gr}'_\bullet(N)\arrow[d,"\operatorname{gr}'(p)"]\\
(\operatorname{gr}(M)\otimes_A(A/\underline J))[T_1,\ldots,T_n]
 \arrow[r,"\psi_M"]\arrow[d]
\&\operatorname{gr}'_\bullet(M)\arrow[d]\\
0\&0\\
|[alias=callout,draw,ellipse]|\text{donner d'autres noms!}\&
\end{tikzcd}}
où les colonnes sont exactes.

\uline{Proposition\deleted{s}} (19.6.3) .- \uline{Avec les notations précédentes, supposons \inserted{que} la
suite} $f$ \inserted{soit} $\operatorname{gr}_\bullet(N)$-\uline{régulière. Considérons les conditions suivantes} :

a) $f$ \uline{est} $\operatorname{gr}_\bullet(M)$-\uline{régulière}.

b) $\psi_M$ \uline{est bijectif}.

c) $\psi'_R$ \uline{est surjectif} (autrement dit $\operatorname{gr}'_\bullet(R)$ est engendré comme
$\operatorname{gr}_{\underline J}^\bullet(A)$-module par
$\operatorname{Im}(\operatorname{gr}_\bullet(R)\longrightarrow\operatorname{gr}'_\bullet(R))$).

d) $\psi'_R$ \uline{est injectif}.

d') \uline{L'homomorphisme} $(\psi'_R)_0:
\operatorname{gr}_\bullet(R)\otimes_A(A/\underline J)\longrightarrow\operatorname{gr}'_\bullet(R)$
\uline{obtenu en restreignant} $\psi'_R$ \uline{aux polynômes de degré} $0$, \uline{est injectif}.

\deleted{\unsure{donc}} e) $\psi'_R$ \uline{est bijectif}.

f) \figure{\tikz[baseline=(a.base)]{\node[draw,ellipse](a){$\operatorname{gr}(j)\otimes1$};}}
\uline{est injectif}.

f') \uline{L'homomorphisme}
$\unsure{\operatorname{gr}_\bullet|\operatorname{gr}^\bullet_\bullet}(R)\otimes_A(A/\underline J)
 \longrightarrow\operatorname{gr}_\bullet(N)\otimes_A(A/\underline J)$
\uline{obtenu en restreignant}
\figure{\tikz[baseline=(a.base)]{\node[draw,ellipse](a){$\operatorname{gr}(j)\otimes1$};}}
\uline{aux polynômes de degré} $0$, \uline{est injectif}.

g) \uline{La filtration} $\underline J$-\uline{préadique de}
$\operatorname{gr}_\bullet(R)$ \uline{est induite par celle de}
$\operatorname{gr}_\bullet(N)$.

h) \uline{Les filtrations} $(R'_k)$ \uline{et} $(R''_k)$ \uline{sur} $R$ \uline{sont identiques}.
\mspage{86}
